\documentclass[10pt,a4paper]{article}
\usepackage[a4paper,margin=18mm,headheight=14pt]{geometry}
\usepackage{fontspec}
\usepackage{amsmath}
\usepackage{amssymb}
\usepackage{xcolor}
\usepackage{booktabs}
\usepackage{tabularx}
\usepackage{enumitem}
\usepackage{listings}
\usepackage{xurl}
\usepackage{hyperref}
\usepackage{fancyhdr}

\setmainfont{Noto Sans CJK KR}
\setsansfont{Noto Sans CJK KR}
\setmonofont{D2Coding}
\hypersetup{
  colorlinks=true,
  linkcolor=blue!45!black,
  urlcolor=blue!45!black,
  pdftitle={문제 6 Dual EC DRBG 백도어 복원과 최적화},
  pdfauthor={암호분석경진대회 제출 보고서}
}
\urlstyle{same}
\setlist{nosep,leftmargin=5mm}
\setlength{\parindent}{0pt}
\setlength{\parskip}{4pt}
\setlength{\emergencystretch}{3em}
\pagestyle{fancy}
\fancyhf{}
\lhead{2026 암호분석경진대회 문제 6}
\rhead{Dual\_EC\_DRBG 분석 및 최적화}
\cfoot{\thepage}

\definecolor{codebg}{HTML}{F4F6F8}
\definecolor{accent}{HTML}{174A7E}
\definecolor{note}{HTML}{F7F3E8}
\lstset{
  basicstyle=\ttfamily\small,
  backgroundcolor=\color{codebg},
  frame=single,
  rulecolor=\color{black!15},
  breaklines=true,
  columns=fullflexible,
  keepspaces=true,
  showstringspaces=false,
  aboveskip=5pt,
  belowskip=5pt
}

\newcommand{\resultbox}[1]{%
  \begin{center}
  \fcolorbox{accent}{blue!4}{\parbox{0.92\linewidth}{\centering
  \large\bfseries\ttfamily #1}}
  \end{center}
}
\newcommand{\notebox}[1]{%
  \begin{center}
  \fcolorbox{orange!55!black}{note}{\parbox{0.92\linewidth}{#1}}
  \end{center}
}

\begin{document}

\begin{center}
  {\LARGE\bfseries 문제 6 Dual\_EC\_DRBG}\\[2mm]
  {\large telemetry 백도어 복원, 다음 출력 예측과 구현 최적화}\\[3mm]
\end{center}

\resultbox{r3 = 0x2443c8daf1a9d52b09}

\begin{tabularx}{\linewidth}{@{}lX@{}}
\toprule
복원량 & 값\\
\midrule
비밀 스칼라 \(d\) &
\texttt{0x1c3cdd6b221806db0a7b28}\\
\(r_0\) 경로에서 얻는 \(s_2\) &
\texttt{0x638d9d631ab436da51e640}\\
\(r_1\) 경로에서 직접 얻는 \(s_3\) &
\texttt{0x948173253ad6d120a3f562}\\
다음 출력 \(r_3\) &
\texttt{0x2443c8daf1a9d52b09}\\
\bottomrule
\end{tabularx}

\section{공격 개요}

문제의 생성식은 내부 상태 \(s_i\)에 대해
\[
  s_{i+1}=X(s_iP),\qquad
  r_i=\operatorname{TMSB}(X(s_{i+1}Q))
\]
이다. 공개 출력은 88비트 \(x\)좌표의 하위 16비트를 버린 72비트다.
telemetry의 여섯 행은 같은 비밀 스칼라 \(d\)에 대한 affine residue의
하위 20비트를 버린 값이다. 풀이의 전체 흐름은 다음과 같다.

\begin{enumerate}
  \item telemetry 생성식을 추론하고 여섯 행으로 \(d\)를 유일하게 복원한다.
  \item 독립 점 등식 \(P=dQ\)로 추론한 식과 \(d\)를 검증한다.
  \item 공개 출력 하나에 하위 16비트를 붙여 가능한 곡선점을 lift한다.
  \item 숨은 관계로 다음 내부 상태를 얻고 나머지 공개 출력을 재생해
        잘못된 후보를 제거한다.
  \item 유일한 상태에서 \(r_3\)를 계산한다.
\end{enumerate}

\notebox{\textbf{상태 표기.}
\(r_0\)을 lift한 점은 \(s_1Q\)이고 \(X(d(s_1Q))=X(s_1P)=s_2\)다.
같은 공격을 \(r_1\)에 적용하면 직접 복원되는 값은 \(s_3\)다.
제출 답과 모든 검증에서 이 index를 구분했다.}

\section{telemetry에서 비밀 스칼라 복원}

\subsection{누설식 추론과 검증}

부분군 위수
\[
  n=\texttt{0x2b674bdfd6fc4ba4ba751d},\qquad B=2^{20}
\]
에 대해 각 행의 \texttt{scale}, \texttt{offset}, \texttt{summary}를
다음의 가장 단순한 affine-residue 모형으로 해석했다.
\[
\begin{aligned}
u_i&=(\operatorname{scale}_i d+\operatorname{offset}_i)\bmod n,\\
\operatorname{summary}_i&=\left\lfloor u_i/B\right\rfloor.
\end{aligned}
\]
이는 문제에서 식으로 주어진 사실이 아니라 검증해야 하는 가설이다.
여섯 \(\operatorname{scale}_i\)가 모두 \(n\)과 서로소이므로 역원이 존재한다.
첫 행에 대해
\[
 (\operatorname{scale}_0d+\operatorname{offset}_0)\bmod n
   =\operatorname{summary}_0B+\ell,\qquad 0\le\ell<B
\]
이고, \(v=\operatorname{scale}_0^{-1}\bmod n\), \(\ell=0\)일 때의 후보를
\(d_0\)라 하면
\[
 d(\ell)=d_0+\ell v\bmod n
\]
이다. 이 식을 둘째 행에 넣으면 미지수 하나 \(\ell\)에 대한 modular
interval 문제가 된다.

\subsection{\texttt{floor\_sum}을 이용한 비열거 복원}

\[
 a=\operatorname{scale}_1v\bmod n,\quad
 c=\operatorname{scale}_1d_0+\operatorname{offset}_1\bmod n
\]
이라 놓으면 둘째 행은
\[
 (a\ell+c)\bmod n\in[L,U)
\]
가 된다. 여기서 \(L=\operatorname{summary}_1B\),
\(U=\min((\operatorname{summary}_1+1)B,n)\)다.

Euclidean recurrence로
\[
 F(k,m,a,b)=\sum_{i=0}^{k-1}
  \left\lfloor\frac{ai+b}{m}\right\rfloor
\]
를 계산하면 다음 식으로 임의의 \(\ell\) 구간에 조건을 만족하는 원소가
몇 개인지 \(O(\log n)\)에 셀 수 있다.
\[
\#\{i<k:(ai+b)\bmod m<y\}
=k-\{F(k,m,a,b+m-y)-F(k,m,a,b)\}.
\]
해가 있는 구간만 이분하여 유일한
\(\ell=\texttt{0x1f051}\)을 찾았다. 복원한
\[
d=\texttt{0x1c3cdd6b221806db0a7b28}
\]
는 여섯 행 모두에서
\[
 ((\operatorname{scale}_i d+\operatorname{offset}_i)\bmod n)\gg20
  =\operatorname{summary}_i
\]
를 만족한다. 또한 전체 좌표를 사용한 독립 검증 \(P=dQ\)도 성립한다.
따라서 알려지지 않은 누설식을 telemetry 내부 일치만으로 순환 검증한 것이
아니다.

기존 \(2^{20}\) 전수조사는 warm 반복 중앙값 1,453.475ms였고,
\texttt{floor\_sum} 경로는 0.750ms였다. 이 인스턴스에서는 약 1,938배
차이이며, 복잡도는 \(O(B)\) modular operation에서
\(O(\log B\log n)\)으로 줄었다.

\section{출력 lift와 \texorpdfstring{\(r_3\)}{r3} 예측}

공개 출력 \(r_j\)마다
\[
 x=(r_j\ll16)\mathbin{|}\ell,\qquad 0\le\ell<2^{16}
\]
를 만들고 \(x<p\)이며
\[
 y^2=x^3+ax+b\pmod p
\]
가 해를 갖는 후보만 남긴다. \(+y\)와 \(-y\)는 서로 반대점이고
\(X(dR)=X(-dR)\)이므로 한 부호만 검사한다.

기본 경로는 \(r_0\)을 lift한 \(R=\pm s_1Q\)에서
\[
 s_2=X(dR)
\]
를 얻는다. 이 상태로 \(r_1\)을 재생해 후보를 거르고,
\(s_3=X(s_2P)\)에서 \(r_2\)도 재생한다. 유일한 생존 후보에 대해
\[
 s_4=X(s_3P),\qquad
 r_3=\operatorname{TMSB}(X(s_4Q))
\]
를 계산하면 제출 답을 얻는다.

최종 native 탐색은 같은 원리를 \(r_1,r_2\) 관측 window에 적용한다.
\(r_1\)의 lift \(T=\pm s_2Q\)에서 \(X(dT)=s_3\)를 얻고 \(r_2\)로
필터한 뒤 \(r_3\)를 계산한다. 정답 low가 \texttt{0x5338}에서
\texttt{0x3cea}로 바뀌어 순차 prefix가 21,305개에서 15,595개로
26.8\% 줄었다. 40개 adjacent AB/BA pair에서 이 변경의 paired
median은 \(1.3428\times\), bootstrap 95\% CI는
\(1.3336\ldots1.3510\)이며 stationarity gate를 통과했다.

\section{최종 구현}

\subsection{고정폭 체와 타원곡선 연산}

최고 성능 C++20 경로는 88비트 소수체를 16바이트 2-limb Montgomery
residue로 표현한다. BMI2/ADX target에서는 \texttt{\_mulx\_u64}와
carry/borrow intrinsic을 쓰고, 그 외에는 portable
\texttt{unsigned \_\_int128} 경로로 fallback한다. Jacobian point는
48바이트 POD다.

임의점 \(d\)배는 원곡선을 \(a=-3\)인 동형 곡선으로 옮겨 Hamburg
co-Z \(x\)-only ladder를 사용한다. exceptional denominator는
width-2 NAF로 fail-safe 처리한다. 정상 경로에는 \(y\)가 필요 없으므로
모든 lift에서 제곱근을 계산하지 않고 Montgomery residue에 직접 적용한
hybrid 128/64비트 Euclidean Jacobi symbol로 제곱잉여를 먼저 판정한다.
실제 제곱근은 NAF fallback이 필요한 예외에서만 지연 계산한다.

반복되는 고정점 \(Q\) 곱에는
\[
 \operatorname{table}[i][j]=j\,2^{8i}Q
 \quad(0\le i<11,\ 0\le j<256)
\]
인 byte-comb를 사용한다. table을 한 번 batch-normalize해
90,112바이트 affine 배열로 저장하므로 이후 고정점 곱은 최대 11회의
mixed addition으로 끝난다.

\subsection{cofactor-5 부분군 필터}

이 곡선은 \(\#E(\mathbb F_p)=5n\), \(\operatorname{ord}(Q)=n\)이다.
따라서 curve-valid lift의 약 \(1/5\)만 올바른 부분군에 있다. 모든 후보에
비싼 Hamburg 곱을 하기 전에 \(x\)-only Frobenius--Tate trace로
\([n]T=O\)와 동치인 부분군 membership을 판정한다.

\(p\bmod5=4\)인 dual embedding degree 2를
\(\mathbb F_{p^2}\) Frobenius \(-1\) eigenspace의 order-5 점에
특수화했다. Miller 값에서
\[
 W=f(T)^p/f(T),\qquad \tau=W+W^{-1}
\]
를 만들면 \(y\)가 소거되어 \(x\)만으로 계산할 수 있다.
\(E=(p+1)/5=20H\)라 두고 Lucas recurrence로 \(L_H\)를 계산한 뒤
\(\mu_{20}/\{z\sim z^{-1}\}\)의 11개 trace와 비교한다.
고정 PRAC schedule은 115 byte, \(118M+6S=124\) products로,
binary Lucas의 \(85M+85S=170\)보다 27.1\% 적다.

정답 prefix의 curve-valid lift 7,713개에서 직접 \([n]T=O\)인 후보
1,547개와 trace 판정 1,547개가 전부 일치했다. 이 필터는 Hamburg 호출의
79.94\%를 제거한다.

\subsection{shifted-square trace와 batch inversion}

expanded trace의 공통 인수를 제거하고 \(z=(x-\alpha)^{-1}\)로 바꾸면
처음에는 5차 reciprocal polynomial을 Horner \(5M\)으로 계산할 수 있다.
다섯 계수는 다시 정확히
\[
\begin{aligned}
r&=\lambda z,\\
\tau&=2+r((r+h)^2+k)^2
\end{aligned}
\]
로 인수분해된다. \(\lambda^5\)는 최고차항 계수다. 차수 5를 만드는 데
field product 두 개만으로는 부족하므로 \(2S+1M\)은
multiplicative-degree 하한을 달성한다.

binary-GCD 역원이 마지막에 Montgomery domain으로 들어갈 때
\(\lambda R^2\)를 곱해 \(r=\lambda z\)를 추가 product 없이 얻었다.
또한 Montgomery batch inversion의 첫 forward와 마지막 두 reverse
endpoint product를 생략했다. 활성 원소가 \(m\)개일 때
normalization+trace 비용은
\[
  8m+I\quad\longrightarrow\quad 6m-3+I
\]
로 줄고, PRAC까지 포함한 block membership은
\[
  130m-3+I
\]
이다. 3-thread 이상 scalar 경로는 후보당 \(127M/S+I\)다.
\(I\)는 field inversion이며 block에서는 하나를 공유한다.

\subsection{연속 cubic recurrence와 병렬 정책}

연속 low-bit 후보는 고정 step \(h\)에 대해
\(f(x)=x^3+ax+b\)를 계산한다. 각 block 첫 원소에서만 직접 평가하고
\[
\begin{aligned}
\Delta f   &=3x^2h+3xh^2+h^3+ah,\\
\Delta^2 f &=6xh^2+6h^3,\\
\Delta^3 f &=6h^3
\end{aligned}
\]
를 field addition으로 갱신한다. recurrence가 곡선 우변을 공급하므로
Hamburg용 \(x^2\)는 Jacobi 직후가 아니라 부분군 필터를 통과한 후보만
계산한다. 정답 prefix에서 field square 6,166회를 제거한다.

작업은 atomic counter로 연속 candidate block을 배분한다. 검증된 adaptive
정책은 1 thread에서 block 64, 2 threads에서 block 32, 3 threads
이상에서 scalar chunk 64다. 2-thread 정책만 별도의 고정 CPU
40-pair 측정에서 승격됐으므로 다른 thread 수로 성능을 보간하지 않았다.
AVX2에는 직접적인 packed \(64\times64\to128\) 정수 곱이 없어 radix
분해와 lane compaction 비용이 커지는 까닭에 기본 경로로 채택하지 않았다.

\section{정확성 검증}

성능 후보는 timing 전에 다음 oracle을 통과해야 한다.

\begin{itemize}
  \item canonical arithmetic와 2,000개 field vector, 64개 경계 pair
  \item affine reference와 256개 point/scalar vector
  \item 실제 lift 128개에서 Hamburg/NAF 및 scalar/batched subgroup 대조
  \item Sage 무작위 200점과 실제 prefix 전체에서 trace와 직접
        \([n]T=O\) 비교
  \item expanded, reciprocal, Horner, shifted-square 식의 경계·무작위·
        공개 prefix \(3\cdot65{,}536\)개 cross-product identity
  \item scaled binary inverse와 일반 inverse oracle, endpoint-elided
        batch의 경계 크기·분모 0 fail-closed 검사
  \item \(P=dQ\), \(d\), 상태 label, 정답 low bits와 \(r_3\) known answer
\end{itemize}

최종 native \texttt{--self-test --json}은 다음을 출력했다.

\begin{lstlisting}
{"self_test":true,
 "field_vectors":2000,
 "field_boundary_pairs":64,
 "point_vectors":256,
 "hamburg_lift_vectors":128,
 "subgroup_lift_vectors":128}
\end{lstlisting}

실제 공격 결과도 다음 known-answer 필드를 모두 확인한다.

\begin{lstlisting}
d=0x1c3cdd6b221806db0a7b28
state_label=s3
state=0x948173253ad6d120a3f562
lift_low_bits=15594
r3=0x2443c8daf1a9d52b09
p_equals_dq=true
\end{lstlisting}

\section{반복 측정과 해석 한계}

측정 환경은 AMD EPYC 7B12 VM, 8 logical CPU, Python 3.11.2,
G++ 12.2.0이었다. 큰 단계 비교는 1회 warm-up 뒤 완전한 공격을 5회
교차 실행했다. 작은 후보는 adjacent AB/BA 40쌍, bootstrap 95\% CI,
실행 순서 stratum과 시간 block stationarity를 함께 검사했다. 매 timing
전에 정답과 실제 thread/schedule/backend metadata를 확인했다.

\begin{tabularx}{\linewidth}{@{}Xrrr@{}}
\toprule
역사적 구현 & 중앙값 & Python 대비 & paired\\
\midrule
기존 Python & 14.298741s & \(1.00\times\) & \(1.00\times\)\\
최적화 Python \texttt{int} & 3.272242s & \(4.37\times\) & \(4.41\times\)\\
최적화 Python \texttt{gmpy2} & 3.000356s & \(4.77\times\) & \(4.61\times\)\\
C++/GMP 1 thread & 1.873220s & \(7.63\times\) & \(7.44\times\)\\
C++/GMP/OpenMP 8 threads & 0.447919s & \(31.92\times\) & \(30.75\times\)\\
\bottomrule
\end{tabularx}

위 표는 언어/backend 전환 당시의 동일 campaign 결과이며, 후속 Jacobi,
부분군 trace/PRAC, shifted-square와 adaptive 정책까지 포함한 현재
source의 절대 시간표가 아니다. 서로 다른 host 부하의 시간을 나눠 더 큰
속도 향상으로 소급하지 않았다.

안정적으로 stationarity gate를 통과한 구성요소 효과는
\(r_1,r_2\) window \(1.3428\times\),
Hamburg \(1.1716\times\), sqrt 대신 Jacobi \(1.0819\times\),
2-thread adaptive block \(1.2121\times\)였다.

shifted-square/Horner는 paired \(1.0077\times\)
(CI \(1.0043\ldots1.0149\))였지만 2\% 승격 문턱과 절대-time
stationarity를 통과하지 못했다. 다만 분리한 normalization+trace cycle은
활성 원소 32/128/256개에서
\(1.269/1.304/1.313\times\)였고,
\(8m+I\to6m-3+I\)의 정확한 operation 감소와 Sage/C++ 등가 검증이
성립하므로 기본값으로 채택했다.

cubic recurrence와 post-subgroup \(x^2\) 지연도 wall-clock
stationarity를 통과하지 못했다. 이들은 각각 direct cubic multiplication과
실제로 버려지는 후보의 square를 제거하는 진부분집합 최적화이며,
전수 등가 검사와 기존 경로를 재현하는 ablation macro를 함께 남겼다.
따라서 이 보고서에서 ``stationarity 실패''는 정확성 실패가 아니라 해당
공유 VM 측정으로 작은 성능 차이를 확정하지 못했다는 뜻이다.

\section{복잡도와 메모리}

\begin{itemize}
  \item telemetry 복원:
        \(O(\log B\log n)\), \(B=2^{20}\).
  \item 상태 복원:
        최악 \(2^{16}\)개의 \(x\) 후보에 Jacobi와 부분군 판정을 수행하고,
        생존 후보에 고정 \(d\) 곱과 공개 출력 검증을 수행하므로
        \(O(2^{16}\log n)\).
  \item 다음 출력:
        상태가 정해진 뒤 상수 횟수의 scalar multiplication.
  \item 고정 메모리:
        90,112바이트 affine \(Q\) table.
        그 밖의 탐색 상태는 worker마다 상수 크기이며 전체 \(2^{16}\)
        후보를 저장하지 않는다.
\end{itemize}

\section{재현}

저장소 루트에서 다음처럼 최종 native source를 빌드한다.

\begin{lstlisting}
g++ -O3 -DNDEBUG -march=native -std=c++20 -fopenmp \
  solutions/06_optimization/deep_native_06.cpp \
  -o /tmp/deep_native_06

/tmp/deep_native_06 --self-test --json
/tmp/deep_native_06 --threads 1 --schedule adaptive --json
\end{lstlisting}

PRAC schedule과 부분군 수학의 독립 감사는 다음 명령으로 재현한다.

\begin{lstlisting}
python3 solutions/06_optimization/generate_06_prac_schedule.py \
  --json
sage -python solutions/06_optimization/audit_06_subgroup.py \
  --json
sage -python solutions/06_optimization/audit_06_subgroup.py \
  --samples 0 --full-prefix --json
\end{lstlisting}

PRAC 생성기는 \texttt{source\_match=true},
\texttt{length=115}, \texttt{field\_products=124}를 출력해야 한다.
Sage full-prefix 감사는
\texttt{prefix\_valid\_lifts=7713},
\texttt{prefix\_members=1547}을 출력해야 한다.

\section{참고 자료}

\begingroup
\raggedright
\begin{enumerate}[leftmargin=6mm]
  \item NIST, \textbf{SP 800-90 Rev. 1} (withdrawn).
        \url{https://csrc.nist.gov/pubs/sp/800/90/r1/final}
  \item D. Shumow and N. Ferguson,
        \textbf{On the Possibility of a Back Door in the NIST SP800-90
        Dual EC PRNG}.
        \url{https://rump2007.cr.yp.to/15-shumow.pdf}
  \item AtCoder Library, \texttt{floor\_sum} documentation and
        implementation.
        \url{https://atcoder.github.io/ac-library/production/document_en/math.html}
  \item Explicit-Formulas Database,
        short-Weierstrass Jacobian formulas.
        \url{https://www.hyperelliptic.org/EFD/g1p/auto-shortw-jacobian.html}
  \item M. Hamburg,
        \textbf{Faster Montgomery and double-add ladders for short
        Weierstrass curves}, ePrint 2020/437.
        \url{https://eprint.iacr.org/2020/437}
  \item N. Möller, \textbf{Efficient computation of the Jacobi symbol}.
        \url{https://arxiv.org/abs/1907.07795}
  \item A. Koshelev,
        \textbf{Subgroup membership testing on elliptic curves via the
        Tate pairing}, ePrint 2022/037 및 출판 correction.
        \url{https://eprint.iacr.org/2022/037.pdf}
  \item A. Enge, \textbf{Bilinear pairings on elliptic curves}.
        \url{https://arxiv.org/abs/1301.5520}
  \item P. L. Montgomery,
        \textbf{Evaluating recurrences of form
        \(X_{m+n}=f(X_m,X_n,X_{m-n})\) via Lucas chains}.
        \url{https://cr.yp.to/bib/1992/montgomery-lucas.pdf}
  \item P. L. Montgomery,
        \textbf{Modular Multiplication Without Trial Division}.
        \url{https://doi.org/10.1090/S0025-5718-1985-0777282-X}
  \item T. Mytkowicz et al.,
        \textbf{Producing Wrong Data Without Doing Anything Obviously Wrong!}
        \url{https://sape.inf.usi.ch/publications/asplos09.html}
\end{enumerate}
\endgroup

\end{document}
